\documentclass[11pt]{letter}
\usepackage[margin=1in]{geometry}
\usepackage[T1]{fontenc}
\setlength{\parskip}{0.55em}
\setlength{\parindent}{0pt}
\signature{Ryota Matsuki\\Independent Researcher, Matsuyama, Ehime, Japan\\ryota.matsuki@gmail.com}
\address{}
\begin{document}
\begin{letter}{Editor-in-Chief\\Journal of Industry, Competition and Trade}
\opening{Dear Professor Haucap,}

Please consider the manuscript \emph{Private Compatibility and the Stability of Standards Coalitions} for publication in the \emph{Journal of Industry, Competition and Trade}.

The paper studies a policy interaction between private technological adaptation and formal standards coordination. Governments choose among fragmented, regional, and international standards arrangements; firms can subsequently pay a fixed cost to support an additional standard. In the canonical three-country Cournot microfoundation, re-solving the market after incomplete outsider adaptation yields an equilibrium residual-rent term that separates three cases: private adaptation can reverse a regional member's ranking toward international standardization, strengthen that incentive without reversing the ranking, or weaken the incentive.

The paper does not claim that this political effect is general across market structures. The analysis establishes a nonempty open set with incomplete adaptation and a selection-free private continuation, then reports pre-specified failure tests. The initial regional advantage disappears at the zero-network boundary, under one alternative singleton-network specification, and in a differentiated-demand comparison under both Cournot and Bertrand competition, while one-way private adoption can remain feasible. A full-bypass coalition-stability exercise is presented as an institutional application rather than as the source of generality.

I believe this conditional applied-theory contribution fits the journal's interest in the microeconomic foundations of market functioning, firm strategy, competition, innovation and new technologies, international trade, and policy-relevant applied theory.

The submission includes an anonymous reviewer manuscript, a separate title page, and supplementary replication material. Symbolic and numerical verification code and reproducible figure/table scripts are supplied. No empirical data are used. The manuscript transparently discloses the use of ChatGPT as an AI-assisted research and manuscript-preparation tool; all mathematical, bibliographic, interpretive, and final editorial judgments were reviewed and accepted under the author's responsibility.

The work is not being resubmitted to IJIO. Final originality, simultaneous-submission, and prior-dissemination certifications will be confirmed in the live submission system immediately before submission.

\closing{Sincerely,}
\end{letter}
\end{document}
